
de pont-levis). C'est sans à-coups que les Protestants vont donc occuper Alban jusqu'à l'hiver 1588. L'autorité centrale est donc mise en échec et les sergents royaux n'osent s'aventurer en ces parages par « peur d'être tués sur place ». Seul le village de Saint-André a une garnison catholique. C'est le capitaine Passemar de Saint-André qui la commande et cela, victorieusement. En effet, il met en échec les troupes du comte de Montmorency, gouverneur protestant de Castres, qui ont tenté de s'emparer de la petite cité, en 1587.

\lettrine{A}{u plan ecclésiastique, c'est en cette année 1588 que Notre Dame d'Alban n'a plus de « Prieur ». Son desservant a titre de « Vicaire perpétuel ». Il est présenté par le prieur bénédictin de Castres mais canoniquement institué par l'évêque d'Albi. Le clergé régulier (bénédictin) cède le pas au clergé séculier diocésain.}

1588 est vraiment une année néfaste pour les Albanais. Les hostilités ont repris dès l'hiver. En décembre, le baron d'Ambres, seigneur catholique, s'empare de la place d'Alban pour en déloger les protestants. Par ses pionniers venus d'Albi il fait raser toutes les défenses de la ville. Avec ses troupes, il se rend ensuite à Curvalle, puis à Villeneuve du Tarn où il n'essuie que des échecs. Sur ces entrefaites, le comte de Montgomery est remplacé à Castres par Turenne. Ce dernier envoie un homme aux nouvelles « pour se rendre compte de ce qui se passe au fort d'Alban qui protège la montagne ». Notre cité n'a donc pas été totalement démantelée et le baron d'Ambres peut y loger une garnison de vingt hommes sous les ordres du Capitaine Durand du Mas Nadal. Ce n'est pas formidable comme moyen de défense et, en juillet 1589, ils capitulent devant les soldats huguenots. En août 1590, derechef le baron d'Ambres s'empare d'Alban. Il y tient trois mois garnison, s'efforçant de fortifier la ville. Mais, le 2 septembre 1591, les Protestants redeviennent maîtres de la situation. Il est vrai pour peu de temps, puis, en effet, le baron d'Ambres les déloge une fois de plus quatre jours après. C'est le 6 septembre 1591. La ville est abandonnée quelques jours plus tard avec pour toute défense une garnison de trois hommes... Deux mois après, les protestants que commandent trois capitaines, Caussé, Portal, Petras, arrivent de Castres avec un canon... Les trois défenseurs de la cité capitulent le 11
